% Vista editorial exacta de corpus/source/masas/04_persistencia.tex, líneas 365--590. Las líneas 1--364 ya residen exactamente en 73_rev2_electron_lectores_torres.tex.
\subsection{Persistence of the enriched state under the surviving extension and conservation of the route memory}
Let \(x\) be an APP state, \(\tau\) its tritic word, and \(B_K\) the memory block at level \(K\). The state relevant to mass has the form

\[
X_K=(x_K,\tau_K,\varepsilon_K,h_K,v_K,c_K,B_K,g_K),
\]
where \(\varepsilon_K\) is the sheet orientation, \((h_K,v_K)\) the cellular orientation, \(c_K\) the carry, and \(g_K\) the nonadic phase. An extension survives when it forms a compatible section of the successive cylinders. If \(C_{K,j}\), \(0\le j<729\) are the children of the cylinder at the next level, the pruning rule satisfies

\[
\sum_{j=0}^{728}G_9(C_{K,j})=1.
\]
The unique index \(j_K\) produces the transition

\[
P_{K+1}=729P_K+j_K.
\]
As nine levels are advanced, the phase is preserved,

\[
g(K+9)=g(K),
\]
but they change the block, prefix, and memory. Thus phase return is not state return. A stable particle is defined by the compatibility of the section across these returns, not by the motionless permanence of a visible coordinate.

The observable route \(\gamma_{\rm obs}\) is the chronological shadow of the extension. The point traversing a diagram is a reading cursor; the particle is the complete persistence class. This distinction permits a correct interpretation of quantum motion: what is required is not a small sphere following a classical trajectory, but a coherent history that admits several compatible projections.

\Needspace{7\baselineskip}
\subsection{Dodecaphasic registration and signature}
The temporal wheel contains twelve oriented events. Its causal decomposition is

\[
12=1+5+6:
\]
one event fixes the anchor, five determine the visible incidence, and six form the opposite pairs surviving the oriented quotient. If \(s_C^{\rm tot}\) is the integer contra-angular sum over the twelve events, the record preserves

\[
s_C:=s_C^{\rm tot}\in\mathbb Z,
\qquad
s_C^{(12)}:=\frac{s_C^{\rm tot}}6.
\]
The first object is the integer coordinate of the signature; the second is its dodecaphase angular reading. The divisor is not a normalization introduced from an external unit. It is the index that appears when the six opposite pairs are reconstructed. The dodecaphase incidence matrix has Smith normal form with factors \((1,12)\); the character is well defined on the quotient because the orientation residue is measured in that structure.

This is the governing convention of the entire work. The preserved historical expressions that designated the already divided coordinate as \(s_C\) are read as \(s_C^{(12)}\); their numerical values and attestations remain intact, but no current formula divides the same coordinate twice.

The direct winding pair is

\[
W_+=6n_A+s_C^{\rm tot},
\qquad
W_-=6n_A-s_C^{\rm tot}.
\]
As \(s_C=s_C^{\rm tot}\) in the published signature, this formula coincides with \(W_\pm=6n_A\pm s_C\). The contra-angular exponent uses \(s_C^{(12)}=s_C/6\) exactly once; it does not again divide a coordinate that had already been normalized.

The complete record publishes the signature

\[
\sigma(\Gamma)=
(n_A,s_C,k_{\Delta_4},b_{120},b_\zeta)\in\mathbb Z^5,
\qquad
\zeta_{270}=135\Delta_4^2.
\]
The five components have distinct types. \(n_A\) counts the direct angular development; \(s_C\) records the integer counter-angular sum; \(s_C^{(12)}=s_C/6\) is its angular reading; \(k_{\Delta_4}\) reads the global refinement; \(b_{120}\) marks the harmonic channel of height 120; \(b_\zeta\) belongs to the quadratic corrector of height 270. In particular, \(s_{270}\) and \(\zeta_{270}\) are not the same object: the former is an odd trace of the bilayer operator; the latter is a quadratic coordinate of the register.

The word electronic

\[
\tau_{12}=+++---+++---
\]
has linear sum zero, but its par signature is not zero:

\[
(Q_3,Q_4,Q_6,\rho_0,w)=(-12,-4,12,-6,12).
\]
Nor should it be identified with the raw signature of the electronic cycle
\((24,0,12,0,-12)\) or with the affine origin \(v_e=0\). These are four
different levels of reading: word, even signature, raw signature, and
origin of the relative character.

\Needspace{7\baselineskip}
\subsection{Angle, counter-angle, and character}
To make the construction autonomous, fix here the three internal coordinates that enter throughout the law:

\[
\alpha_{\rm HMT}
=\frac{7297352569283800997285105472380663}{10^{36}},
\]
\[
\Delta_4
=\frac{\pi_{\rm HMT}-e_{\rm HMT}\log\pi_{\rm HMT}}{270}
\left(1+\frac{\pi_{\rm HMT}}{729}\right),
\]
and, for \(a>0\) and \(\phi>0\),

\[
H_5(a,\phi)
=\frac12\sqrt{\frac{1000a}{\phi}}-a
+\frac{9a^2}{16}-\frac{5a^3}{9}
+\frac{7a^4}{48}-\frac{a^5}{54}.
\]
Henceforth,

\[
H_5:=H_5(\alpha_{\rm HMT},\varphi_{\rm HMT}).
\]
These definitions do not import an external metrological constant: \(\pi_{\rm HMT},e_{\rm HMT},\varphi_{\rm HMT}\) are the coinductive evaluations of the HMT reader, and \(\alpha_{\rm HMT}\) is the rational output of dodecaphase closure. The internal constant \(\alpha_{\rm HMT}\) precedes the counter-angle. The regional branch first determines

\[
A=1000\alpha_{\rm HMT},
\qquad
D_A=\exp\!\left(-\frac{100\pi\alpha_{\rm HMT}}9\right),
\]
\[
\mathcal C_\pi(\alpha_{\rm HMT})
=169\alpha_{\rm HMT}^{6}
+\frac{D_A\alpha_{\rm HMT}^{7}}
       {1-D_A\alpha_{\rm HMT}},
\qquad
y_*=\frac{\pi}{90}(H_5+\alpha_{\rm HMT})
-\mathcal C_\pi(\alpha_{\rm HMT}),
\]
\[
C^*=\frac{180}{\pi}y_*.
\]
Here \(A\) and \(C^*\) are already constructed angular coordinates, not action magnitudes. If

\[
\alpha_\angle:=\alpha_{\rm HMT}\,{}^\circ,
\qquad
\eta_{\rm ret}^{(\deg)}:=\frac{C^*}{2}-\alpha_\angle,
\]
then

\[
C^*=2\bigl(\eta_{\rm ret}^{(\deg)}+\alpha_\angle\bigr)
\]
is the inverse return identity, not a second rule for selecting the counter-angle. This notation avoids confusing the angular mantissa \(\eta_{\rm ret}^{(\deg)}\) with the dimensional magnitude \(\hbar_{\rm ret}\). The radian charts are obtained by multiplication by \(\pi/180\). The conjugate channels are

\[
q_\pm=\exp\left[-\frac\pi{180}(A\pm C^*)\right].
\]
The central character of a signature is

\[
\chi_0(\sigma)=
\exp\left(
n_AA_{\rm rad}
+\frac{s_C}{6}C^*_{\rm rad}
+k_{\Delta_4}\Delta_4
+b_\zeta\zeta_{270}
\right).
\]
The coordinate \(b_{120}\) remains in the signature, but its contribution is incorporated into the tower family. This convention prevents counting the same height twice.

The character is dimensionless. Its function is not to become an action constant, but to act by positive homotheties on an already typed mass line. For a pair of states \(X,Y\),

\[
\frac{m_Y^{(0)}}{m_X^{(0)}}
=\chi_{\rm full}(\sigma_Y-\sigma_X,\eta_Y-\eta_X).
\]
This formula separates the ratios basal of mass of the choice common of chart absolute. The quotient beta complete incorporates afterward \(R_{\rm act}^{\kappa_Y-\kappa_X}\).

\Needspace{7\baselineskip}
\subsection{Global mass tower semigroup and convergent harmonic expansion}
Let \(R^2=I\), with projectors

\[
P_\pm=\frac{I\pm R}{2},
\qquad
x=\frac{\pi A}{180},
\qquad
y=\frac{\pi C^*}{180},
\]
and

\[
T=e^{-xI-yR}=q_+P_++q_-P_-.
\]
The even and odd traces are

\[
c_n=\operatorname{Tr}(T^n)=q_-^n+q_+^n
=2e^{-nx}\cosh(ny),
\]
\[
s_n=-\operatorname{Tr}(RT^n)=q_-^n-q_+^n
=2e^{-nx}\sinh(ny).
\]
Satisfy

\[
c_n^2-s_n^2=4e^{-2nx},
\]
\[
c_{m+n}=\frac{c_mc_n+s_ms_n}{2},
\qquad
s_{m+n}=\frac{s_mc_n+c_ms_n}{2},
\]
and the second-order linear recurrence

\[
X_{n+2}=c_1X_{n+1}-e^{-2x}X_n,
\qquad X\in\{c,s\},
\qquad c_0=2,\quad s_0=0.
\]
The set of heights is not reduced to eight visible corrections. It is the semigroup

\[
\langle90,120\rangle
=\{90,120\}\cup\{30r:r\ge6\}.
\]
The heights \(90,120,180,210,240,270,360,420\) are the first useful genealogical witnesses, not the entire domain. The refined character is

\[
\log\chi_{\rm full}(\sigma,\eta)
=\log\chi_0(\sigma)
+\sum_{h\in\langle90,120\rangle}N_h(\sigma,\eta)s_h,
\]
with

\[
N_h(\sigma,\eta)=\eta_h+b_{120}\delta_{h,120}.
\]
Thus \(s_{120}\) appears only once, and at height 270, \(N_{270}s_{270}\) and \(b_\zeta\zeta_{270}\) remain separate. Absolute convergence requires

\[
\sum_h |N_hs_h|<\infty.
\]
The completeness of the semigroup guarantees the capacity to represent positive ratios; it does not, by itself, select a species' physical coefficients. The prospective selection rule, when used, must arise from the route, genealogical record, carry, memory, and winding before the external mass is consulted.

\Needspace{7\baselineskip}
